\section{Conclus\~ao e Trabalhos Futuros}
\label{sec:conclusao}

Este trabalho propôs um pipeline que liga as opiniões publicadas sobre audiências públicas à fala de
quem as emitiu, condensa a fala de cada ator em um perfil escrito somente a partir dela e usa esse
perfil para simular como a pessoa se posicionaria, com a base da estimativa declarada e uma
justificação conferida. A UDV fornece evidência localizada para 2.105 das 2.203 opiniões, mas, num
controle com as 183 opiniões de treino que contêm citação confiável, a sentença de maior cosseno
coincide com a passagem citada em 61,2\% dos casos, e a precisão da camada semântica ainda precisa
ser medida pela validação humana. <Síntese dos resultados da simulação.>

A sustentação das falas simuladas precisa ser avaliada por leitura humana, pois a conferência
automática cobre a existência dos itens citados e o início literal das cópias, mas não o sentido.
Ficam em aberto o efeito dos exemplos de fala, que a múltipla escolha não mede, uma condição só com
os trechos recuperados e a comparação entre perfis gerados por modelos e prompts diferentes. A fala
simulada é uma estimativa condicionada às falas públicas da pessoa, e não uma declaração dela: deve
ser apresentada com a base da estimativa e a justificação, nunca como citação, e não serve para
inferir intenção política.
